\documentclass[11pt]{article}
\usepackage[margin=1in]{geometry}
\usepackage[T1]{fontenc}
\usepackage[utf8]{inputenc}
\usepackage{amsmath,amssymb,amsthm,mathrsfs}
\usepackage[hidelinks]{hyperref}
\newtheorem{theorem}{Theorem}
\newtheorem{lemma}[theorem]{Lemma}
\newtheorem{corollary}[theorem]{Corollary}
\newcommand{\DD}{d'd''}
\newcommand{\PH}{\operatorname{PH}}
\newcommand{\MA}{\operatorname{MA}}
\title{Local analytic divisor tests for bounded solution differences\\and the remaining pluriharmonic recognition gate}
\author{Working paper; independent review pending}
\date{9 October 2026}
\begin{document}
\maketitle
\begin{abstract}
On the specified compact analytic torus quotient, equality of bounded
continuous Monge--Ampere measures forces the difference to be harmonic
against every local principal analytic divisor current. We give the
bounded Green argument, globally glued capped-log tests, and the passage
from Laurent polynomials to arbitrary local holomorphic functions by
rational-domain approximation. The last passage uses convergence of
positive currents and continuous multiplication, without assuming that
a continuously weighted closed current is closed. Separately, the
classical pluriharmonic maximum principle gives a precise sufficient
recognition criterion for uniqueness. Recognition of the solution
difference by all curve tests, unbounded restriction, compactness of
finite-energy sublevels, and general equation existence remain open.
\end{abstract}

\section{Carrier, domain, and exact foundations}
Put $K=\mathbb C((s))$, $|s|=e^{-1/8}$, and let
\[
 X=((\mathbb G_m)^2)^{\mathrm{an}}/
 \langle(s^8,2),(3,s^8)\rangle,
 \qquad x_i=-\log|z_i|,
 \quad \rho=\tfrac12|x|^2,
 \quad\theta=\DD\rho.
\]
We retain the compact, connected, smooth, good, boundaryless analytic
quotient and its ordinary injective torus-chart atlas as geometric
foundations. Tropical translation by the periods is $\mathbb Z^2$;
rectangles of coordinate width less than one give injective charts.
The local reference potentials differ by affine functions of unit
logarithms, so their Hessians glue. Quotient existence, separatedness,
properness and nonalgebraization are separate geometric statements;
the present theorem does not reestablish their entire proof trees.

Let $\mathscr P$ be the minimal sheaf of finite continuous functions
containing ordinary smooth CLD psh functions and closed under locally
uniform limits. Let
\[
 \mathcal C=\{u\in C(X,\mathbb R):\rho+u\in\mathscr P
                         \text{ on every chart}\}.
\]
Here every continuous-sheaf stage is retained. Convex combinations and
finite maxima preserve this class: prove the assertions for smooth
weights by convexity and smooth convex soft maxima, then use minimality
and uniform closure successively in the two slots.

The immediately applied CLD v2 foundations are Definition8.1.7 and
Lemma8.1.8; Lemmas8.5.1,8.5.3,8.5.5; the symmetric recursive continuous
products in8.7.1,8.7.2,8.7.4; multilinear continuous dpsh calculus and
smooth null changes in8.8.4--5; scalar Radon recognition and
smooth-coefficient measures in7.2.8,7.2.11--12; compact-support
Green/Stokes in5.11.1; calibrated integration15.7.1; and
Poincare--Lelong16.7.3~\cite{CLD}.
They give positive closed products, local uniform continuity, and the
following multiplier rule: if positive currents $T_j$ converge weakly
to $T$ and continuous coefficients converge locally uniformly, their
products converge weakly. We use only scalar and surface $(1,1)$
positivity; no higher-degree strong-cone order unit is assumed.
Smooth cutoffs and the underlying calibrated integration are ordinary
published foundations. The whole immediate applied source units were
read; deeper recursive proofs are not all newly audited.

We also use the defining Tate-algebra presentation of a rational domain
and the density of finite polynomials in a completed Tate algebra;
the latter follows directly by truncating its convergent series.
The ordinary presentation and its universal-property argument are
Van der Put2.1--2.2~\cite{VdP}; the strict presentation is also stated
in Papikian3.4--3.5~\cite{Pap}.
This includes generalized polyannulus radii. Rational domains form a
compact neighborhood basis inside a good affinoid torus chart.
These are ordinary affinoid foundations, not a claimed approximation
by polynomials on an arbitrary open set without shrinking.

\section{Bounded Green calculus at all continuous stages}
Fix a positive closed $(1,1)$ current $T$. Complete compact smooth
$(1,0)$ forms modulo the null space of the positive pairing
$\int a\wedge Jb\wedge T$. With a nonnegative compact smooth cutoff
$\chi$, write the local gradient pairing as $Q_{T,\chi}$.
Let $\mathscr H_T$ consist of continuous functions locally approximable
uniformly by smooth functions whose gradients converge strongly in
these cutoff seminorms.

Gradients are unique: if smooth $f_j\to0$ uniformly and their gradients
converge strongly to $a$, then, for each compact smooth form $b$,
closedness of $T$ gives
\[
 \int d'f_j\wedge Jb\wedge T
 =-\int f_jd'(Jb)\wedge T\longrightarrow0.
\]
The right side pairs $f_j$ with a finite signed measure. Smooth forms
are dense in the completion, so $a=0$. Smaller patches and cutoffs
give the same local assertion. Smooth partitions patch approximants:
the gradient of $\sum\chi_i f_{i,j}$ splits into
$\sum\chi_i d'f_{i,j}+\sum f_{i,j}d'\chi_i$, with both sums strongly
convergent. Bounded continuous multiplication of forms follows by
scalar smooth approximation and the norm bound. Products and cutoffs
of functions in $\mathscr H_T$ therefore have the usual gradient rule.

If $f\in\mathscr H_T$ and $L_Tf=\DD(fT)$ is signed Radon, smooth
approximation first in $f$, then uniformly and strongly in a compactly
supported $g\in\mathscr H_T$, proves
\[
 \int gL_Tf=-Q_T(f,g).
\]
In particular the product rule yields
\[
 Q_{T,\chi}(w,w)
 =\tfrac12\int w^2\DD\chi\wedge T-\int\chi w L_Tw.
\]
The sheaf of $f\in\mathscr H_T$ with $L_Tf\ge0$ contains smooth psh
functions and is uniformly closed. Indeed choose nested cutoffs
$\chi,\eta$, with $\eta=1$ near $\operatorname{supp}\chi$, and a common
bound $|f_j|\le B$ on the whole support of $\eta$. For
$\epsilon=\|f_j-f_k\|$ there, positivity and integration by parts give
\[
 0\le Q_{T,\chi}(f_j-f_k,f_j-f_k)
 \le\tfrac{\epsilon^2}{2}\|\DD\chi\wedge T\|_{TV}
 +2B\epsilon\|\chi\|_\infty\|\DD\eta\wedge T\|_{TV}.
\]
To obtain smooth approximants to the limit on a smaller patch, choose
an approximant to each $f_j$ with uniform and gradient error at most
$1/j$ for a cutoff equal to one on that patch. The displayed estimate
proves strong convergence. Positivity survives distributional limits.
Minimality now puts all of $\mathscr P$ in this sheaf. Its vector span
contains continuous dpsh functions. Finite patching on compact $X$
proves the global formula
\[
 \int f\DD(gT)=\int g\DD(fT)=-Q_T(f,g),
 \qquad Q_T(f,f)\ge0,
\]
and Cauchy--Schwarz for continuous dpsh functions. This argument uses
positive Laplacians and signed measures, not arbitrary continuous
multiplier closedness.

For $u\in\mathcal C$ write $\theta_u=\theta+\DD u$ and
$B(u,v)=\theta_u\wedge\theta_v$ for the actual recursive local
products, glued through the smooth null potential changes.
Set $\MA(u)=B(u,u)$. Every mixed measure has the same finite mass
$M=\int\theta^2=2$; Green/Stokes annihilates the exact terms.

\section{An elementary bounded all-test estimate}
Let $a,b,w\in\mathcal C$, set $f=a-b$, $g=(a+b)/2$, $h=w-g$, and
\[
 I(a,b)=\int(a-b)(\MA(b)-\MA(a)).
\]
The bounded product expansion gives
\[
 I=Q_{\theta_a}(f,f)+Q_{\theta_b}(f,f)\ge0.
\]
Multilinearity and Green give the signed identity
\[
 Q_{\theta_w}(f,f)-Q_{\theta_g}(f,f)
 =Q_{\theta_a}(f,h)-Q_{\theta_b}(f,h).
\]
Indeed both sides are $-\int f\DD f\wedge\DD h$.
Thus Cauchy--Schwarz implies
\[
 Q_{\theta_w}(f,f)\le I/2+
 \sqrt{I\bigl(Q_{\theta_a}(h,h)+Q_{\theta_b}(h,h)\bigr)}.
\]
For an explicit bounded version choose $B_0$ with
$|a|,|b|,|w|\le B_0$. Then $|h|\le2B_0$, and
\[
 Q_{\theta_a}(h,h)
 =-\int h\,[B(a,w)-(\MA(a)+B(a,b))/2]
 \le4MB_0.
\]
The same holds in the $b$ slot, since every positive measure in the
brackets has mass $M$. Consequently
\[
 Q_{\theta_w}(a-b,a-b)\le I/2+\sqrt{8MB_0I}.
\]
This estimate is sufficient here; it claims no uniform sup-norm bound
for unbounded generated presentations.

\begin{lemma}\label{annihilate}
If $u,v\in\mathcal C$ satisfy $\MA(u)=\MA(v)$ and $d=u-v$, then
$\DD(d\theta_w)=0$ for every $w\in\mathcal C$.
\end{lemma}
\begin{proof}
Equality of the measures gives $I(u,v)=0$. The estimate gives
$Q_{\theta_w}(d,d)=0$. Cauchy--Schwarz with compact smooth scalar
tests and the Green identity show that every test pairing of
$\DD(d\theta_w)$ vanishes. The difference is bounded and continuous.
\end{proof}

\section{Laurent divisor currents from glued smooth caps}
Take a nonzero Laurent polynomial $A$. At a point of its zero locus
choose an injective chart $|x_i-a_i|<r<1/2$, and put
$\rho_a=|x-a|^2/2$. Choose a smooth convex $\chi$ with
$\chi=0$ on $(-\infty,0]$, $0\le\chi'\le1$, and
$\chi(t)=t-c$ for $t\ge1$. The functions
\[
 \psi_N=\chi(\log|A|+N)
\]
extend smoothly by zero near the zeros of $A$ and are ordinary smooth
psh: away from those zeros $A$ is a unit moment, and at each zero the
cap is locally constant.

Fix $0<\delta<r^2/4$ and choose $\epsilon_N>0$ such that
$\delta+\epsilon_N\sup\psi_N<r^2/2$ on the closed polyannulus.
Then
\[
 w_N=\max(\rho_a,\delta+\epsilon_N\psi_N)-\rho_a
\]
extends by zero to a member of $\mathcal C$. The strict boundary
inequality supplies a collar where it is zero, and smooth soft maxima
give the required locally uniform psh approximation. On the fixed
smaller neighborhood $\rho_a<\delta/2$ one has
$\theta_{w_N}=\epsilon_N\DD\psi_N$. Lemma\ref{annihilate} gives
$\DD(d\DD\psi_N)=0$ there.

CLD16.7.3 applies to the regular holomorphic $A$ on the smooth,
pure-dimensional, boundaryless torus chart:
\[
 \DD[\log|A|]=\delta_{\operatorname{div}A}.
\]
The lower-dimensional support fact used in its immediate proof says
that compact smooth top-degree tests avoid $\{A=0\}$.
On such a test support, $\psi_N-N=\log|A|-c$ for all sufficiently large
$N$. Thus the positive currents $\DD\psi_N$ converge weakly to the
actual divisor current. CLD8.5.1 with the continuous coefficient $d$,
followed by continuous differentiation of currents, proves
\[
 \DD(d\delta_{\operatorname{div}A})=0
\]
on the smaller neighborhood. Covering the divisor proves the local
identity everywhere in every chart; off the divisor it is automatic.
The argument does not assert closedness of $d\delta_{\operatorname{div}A}$.

\section{Rational-domain approximation with controlled denominators}
Here is the neighborhood point needed below. The ordinary topology of
an affinoid spectrum is the initial topology of its absolute-value
functions. Around a point in the interior of a polyannulus, choose
finitely many strict lower and upper bounds defining a neighborhood
contained in the prescribed open set. Move the endpoints slightly
inwards, keeping the point inside. Bounds of the form $|s|^{m/n}$ are
dense among positive real bounds, so every chosen constraint can be
written using an integral power of its function and a scalar in $K$.
Weierstrass upper constraints have denominator one. For the finitely
many positive lower constraints $|h_i|\ge a_i$, use common denominator
$h_0=\prod_i h_i$, numerators $h_0/h_i$, and add numerator one with
the redundant bound $1\le(\prod_i a_i)^{-1}|h_0|$. The added numerator
ensures the unit-ideal condition and excludes every zero of $h_0$;
on its complement the other inequalities are exactly the lower
constraints. Upper constraints multiply their numerators by $h_0$.
This gives a rational domain contained in the chosen neighborhood and
having the point in its ordinary interior. Integral powers realize
the scalar bounds in the ordinary strict presentation when needed.
Thus shrinking to such a compact rational neighborhood is justified
in the ordinary topology, not merely by a rigid admissible cover.

\begin{lemma}\label{rational}
Let $V$ be a rational domain in a compact torus polyannulus $P$, and
$F\in\mathcal O(V)$. There are rational functions
$F_n=A_n/Q_n$, with Laurent-polynomial numerators and denominators,
such that $Q_n$ is a unit on $V$ and $F_n\to F$ uniformly on $V$.
\end{lemma}
\begin{proof}
Write the defining presentation as
\[
 \mathcal O(V)=\mathcal O(P)\langle r_i^{-1}T_i\rangle/
                    (g_0T_i-g_i),
\]
where $g_0,g_1,\ldots,g_m$ generate the unit ideal and the domain is
defined by $|g_i|\le r_i|g_0|$. These inequalities and the unit-ideal
condition imply that $g_0$ has no zero on $V$. It is a unit there and
has a positive lower absolute-value bound by compactness.
Finite polynomials in $T_i$ with coefficients in $\mathcal O(P)$
approximate $F$. Each such expression is a finite sum of products of
the $g_i/g_0$ and finitely many coefficients. Laurent polynomials are
dense in the polyannulus algebra by its completed-series definition.
Approximate these finitely many coefficients and all $g_i$ on $P$.
Choose the approximation $q$ to $g_0$ with
$\|q-g_0\|_V<\inf_V|g_0|$. Then $|q|=|g_0|$ on $V$, and $q$ is a
unit there. Substitution in the fixed finite expression is continuous
in the uniform norm, because inverse denominators have this fixed
positive lower bound. Clearing a finite power of $q$ gives a quotient
$A/q^m$ of Laurent polynomials. Choosing each finite truncation and
its substitutions with error tending to zero proves the assertion.
No uniform bound on the degrees or on the denominator powers is needed.
\end{proof}

\section{All local principal analytic divisor tests}
\begin{theorem}\label{localdiv}
Let $u,v\in\mathcal C$ satisfy $\MA(u)=\MA(v)$, and let $W$ be an
ordinary open subset of a torus chart. For every regular holomorphic
function $F$ on $W$,
\[
 \DD\bigl((u-v)\delta_{\operatorname{div}F}\bigr)=0
                      \quad\text{on }W.
\]
This includes local analytic hypersurfaces not defined by a Laurent
polynomial on the entire torus.
\end{theorem}
\begin{proof}
Shrink around any point to the interior of a compact rational domain
$V\subset W$ in a polyannulus. Lemma\ref{rational} gives
$F_n=A_n/Q_n\to F$ uniformly. After omitting initial terms the numerator
is nonzero. A nonzero Laurent polynomial is regular on the smooth
torus, and $Q_n$ is a unit on $V$. Poincare--Lelong therefore gives the
positive current identity on its interior
\[
 T_n=\DD[\log|F_n|]
       =\delta_{\operatorname{div}A_n}.
\]
The preceding section gives $\DD(dT_n)=0$ there, where $d=u-v$.

For any compact smooth $(2,2)$ test $\alpha$ in this interior, its
support avoids the lower-dimensional zero locus of the regular $F$.
Hence $m_\alpha=\inf_{\operatorname{supp}\alpha}|F|>0$. Uniform
approximation implies $|F_n-F|<m_\alpha$ on that support for all late
$n$. The ultrametric inequality gives $|F_n|=|F|$ there. The currents
$[\log|F_n|]$ therefore agree with $[\log|F|]$ on each such test for
all sufficiently late $n$, and converge weakly. Differentiation and
Poincare--Lelong yield
\[
 T_n\longrightarrow\DD[\log|F|]
                          =\delta_{\operatorname{div}F}.
\]
These are positive currents. CLD8.5.1 now gives
$dT_n\to d\delta_{\operatorname{div}F}$, because $d$ is bounded
continuous. Differentiating and using $\DD(dT_n)=0$ proves the result
on the interior of $V$. The local assertions glue on $W$.
\end{proof}
This passage uses approximation on a rational neighborhood and
invertibility of its denominators. It makes no assertion that an
arbitrary analytic curve map has a local closed hypersurface image,
or that the current identity automatically survives every scalar
extension. The weighting step is not licensed for an unbounded
generated difference without a separate integrability theorem.

\section{An exact sufficient recognition criterion}
Gubler--Rabinoff Definition7.1 and Propositions7.2--7.3, with
Theorems7.4 and7.6, concern $\PH$, the functions whose positive and
negative are both classically psh~\cite{GR}. Classical psh means tests
by all analytic curves after every analytic field extension. These
source results give continuity, local maximum and minimum principles,
and finite Shilov determination. They are existing source theorems,
not new uniqueness theorems for the present Monge--Ampere class.

\begin{corollary}\label{criterion}
If the bounded continuous difference $u-v$ in
Theorem\ref{localdiv} is locally classically pluriharmonic, it is
constant on $X$.
\end{corollary}
\begin{proof}
Sheafiness makes it globally pluriharmonic. Compactness supplies a
global maximum. The maximum set is nonempty and closed by continuity,
and open by the local maximum principle at every point of the
boundaryless space. Connectedness makes it all of $X$.
\end{proof}

For completeness, this source framework also gives a finite-probe
compactness statement only for pluriharmonic defects. Cover a
quasicompact analytic space $Z$ by finitely many affinoid domains
$V_i$, and let $\Gamma$ be the union of their finite Shilov boundaries.
The map $h\mapsto(h(\xi))_{\xi\in\Gamma}$ embeds $\PH(Z)$ in
$\mathbb R^\Gamma$, and
\[
 \|h\|_{C(Z)}=\max_{\xi\in\Gamma}|h(\xi)|.
\]
Its image is a finite-dimensional linear subspace, hence closed.
Every bounded family is relatively compact in the uniform norm.
More precisely, convergence of a net on this finite set to a finite
vector implies uniform convergence to a unique pluriharmonic function:
the vector belongs to the closed image and the displayed norm measures
the convergence. No continuity assumption on a separately proposed
pointwise limit is needed. This is a direct finite-dimensional
consequence of the stated source maximum principles.

\section{Reading boundaries and unchanged original residual}
The whole CLD text is the retained v2 body with SHA256
\texttt{64846cb00926267c6238d1afac61b5eaed0a0e3a6349725738f9b1bf163de085}.
The entire applied Poincare--Lelong statement and immediate proof were
freshly read in literal LF23458--23671. The earlier corrected literal-LF
bounded-calculus units are retained with their complete proof-reading
boundaries. No fresh CLD primary pixels or whole recursive foundation
audit is claimed. The GR v2 whole LF1833--1945, including the complete
applied7.1--7.6 statements and immediate proofs, was freshly read.
The later semistable7.8 statement is not a premise here. Full primary
sources, literal units, pixels and book bodies remain local.
For the rational presentation, the Van der Put source's actual original
PDF pages5--6 (printed J7-04--05), including the complete2.2 immediate
proof, are retained with their separately recorded pixel-reading scope.
Papikian's strict-algebra statement has no immediate proof printed in
the retained notes; it is corroboration, not a fresh complete-proof
claim or a theorem about the ordinary Berkovich topology.

All mathematical claims in this paper are UNREVIEWED. In particular,
Corollary\ref{criterion} is a sufficient criterion: local analytic
divisor-current annihilation has not here been proved equivalent to
classical pluriharmonicity by all curves and all extensions. Unbounded
restriction, null-set removability and pointwise recognition remain
missing. The compactness statement for $\PH$ does not establish
compactness or representability of finite-energy psh sublevels.
Intrinsic recognition, actual obstacle envelopes and orthogonality,
existence for every eligible arbitrary-support measure, and full
uniqueness modulo constants remain unresolved for original P16.7.
The present work was not supplied to the independently running
native278 or279 campaigns. Their frozen inputs remain unchanged.

\begin{thebibliography}{9}
\bibitem{CLD} A. Chambert-Loir and A. Ducros,
\emph{Formes differentielles reelles et courants sur les espaces de
Berkovich}, arXiv:1204.6277v2, exact retained source snapshot.
\bibitem{GR} W. Gubler and J. Rabinoff,
\emph{Classically psh and pluriharmonic functions on Berkovich spaces},
arXiv:2506.13548v2, 17 September 2025.
\bibitem{VdP} M. van der Put, \emph{Rigid analytic spaces},
Groupe de travail d'analyse ultrametrique3(2), exp.J7 (1975--1976),
Sections1.1 and2.1--2.2, exact retained Numdam scan.
\bibitem{Pap} M. Papikian, \emph{Rigid analytic geometry and abelian
varieties}, notes dated13 June2004, Sections3.1--3.2, Definition3.4
and Theorem3.5.
\end{thebibliography}
\end{document}
